\section{Method}

\subsection{Overview}

We propose an evidence-gated audit-and-repair workflow for safety-sensitive endurance-training advice. The method is specified as a set of reproducible control points rather than as a system-flow narrative: each control point receives an explicit input, applies a fixed decision rule, writes an output state, and records trace fields for later evaluation.

Figure~\ref{fig:workflow} summarizes the workflow using the module names used throughout this section: Pre-Gate Request Filter, Evidence Gate, Risk Gate, Evidence-Constrained Generator, Independent Grounding Auditor, Bounded Citation/Grounding Repair, Refusal, and Final Answer. A request may end in Refusal only at four refusal-capable control points: the Pre-Gate Request Filter, Evidence Gate, Risk Gate, or Independent Grounding Auditor. The Generator and Repair modules do not independently refuse; they either produce a constrained draft or a repair result that is re-checked by the Independent Grounding Auditor.

\begin{figure}[t]
  \centering
  \includegraphics[width=\linewidth]{figure1_method_workflow.png}
  \caption{Evidence-gated audit-and-repair workflow for safety-sensitive endurance-training advice. User requests first pass through the Pre-Gate Request Filter and evidence retrieval. The Evidence Gate and Risk Gate determine whether generation is allowed, bounded, or refused. Draft answers must pass the Independent Grounding Auditor; Bounded Citation/Grounding Repair handles \texttt{citation repair} for invalid citations, grounding mismatches, unsupported wording, or missing refusal-boundary issues before re-audit. Refusal is an explicit workflow state rather than a generic failure.}
  \label{fig:workflow}
\end{figure}

The workflow is evaluated as S3, the full system. Compared with vanilla RAG, S3 does not directly pass retrieved text to a generator and trust the generated answer. Retrieved or curated evidence is first checked by the Evidence Gate; safety-sensitive cases are then checked by the Risk Gate; draft answers are reviewed by the Independent Grounding Auditor before release as a Final Answer.

\subsection{Output States}

Given a query \(q\), the system returns one of three output states:

\begin{itemize}[leftmargin=*]
  \item \texttt{answered}: a supported answer passes the Independent Grounding Auditor;
  \item \texttt{partial\_answer}: a bounded answer passes audit after \texttt{citation repair}, usually because support constraints required a conservative rewrite;
  \item \texttt{refused}: one of the four refusal-capable control points determines that the request is blocked, insufficiently evidenced, unsafe, or unrepairably unsupported.
\end{itemize}

This design treats refusal as a first-class diagnostic result. In a safety-sensitive advisory setting, a correct refusal can be the intended behavior rather than a model failure.

\subsection{Evidence Modes}

The workflow supports three evidence modes:

\begin{itemize}[leftmargin=*]
  \item \texttt{retrieval\_only}: uses ordinary vector-retrieved contexts from the current knowledge base;
  \item \texttt{gold\_only}: uses curated benchmark evidence spans;
  \item \texttt{retrieval\_plus\_gold}: combines retrieved contexts with curated evidence.
\end{itemize}

These modes separate evidence availability from downstream workflow behavior. \texttt{gold\_only} serves as a curated-evidence upper bound, while \texttt{retrieval\_only} exposes failures caused by retrieval noise or evidence misses. The main 100-question run uses \texttt{retrieval\_plus\_gold}, which is the most informative setting for testing the full workflow while retaining curated benchmark traceability.

\subsection{Reproducible Control Points}

Table~\ref{tab:method-control-points} specifies the four refusal-capable control points. Each row can be implemented deterministically or by a constrained classifier, but the recorded output state and trace fields are fixed for evaluation.

\begin{table}[t]
\centering
\tiny
\setlength{\tabcolsep}{1.5pt}
\caption{Reproducible specification of refusal-capable control points.}
\label{tab:method-control-points}
\begin{tabular}{@{}L{0.15\linewidth}L{0.16\linewidth}L{0.30\linewidth}L{0.18\linewidth}L{0.14\linewidth}@{}}
\toprule
Control point & Input & Decision rule & Output state & Trace field \\
\midrule
Pre-Gate Request Filter & Raw query \(q\). & Block targeted injection, fabrication pressure, unsupported guarantees, red-flag continuation, worsening pain, or safety-suppression patterns. & \texttt{allowed} or \texttt{blocked}; blocked maps to \texttt{refused}. & \texttt{pre\_gate} \\
Evidence Gate & Query, retrieved evidence, optional gold evidence. & Direct support \(\rightarrow\) \texttt{answerable}; bounded support \(\rightarrow\) \texttt{partial}; absent evidence, unsupported precision, or fabricated citation requests \(\rightarrow\) \texttt{unanswerable}. & \texttt{answerable}, \texttt{partial}, or \texttt{unanswerable}; \texttt{unanswerable} maps to \texttt{refused}. & \texttt{evidence\_gate} \\
Risk Gate & Query, evidence state, evidence bundle, risk triggers. & Red flags, illness/heat-illness cues, persistent pain, high fatigue, or unsafe escalation require de-escalation or refusal; continuation, escalation, and diagnosis-like reassurance are forbidden when risk evidence is present. & \texttt{safe}, \texttt{bounded}, or \texttt{unsafe}; \texttt{unsafe} maps to \texttt{refused}. & \texttt{risk\_gate} \\
Independent Grounding Auditor & Draft, cited evidence IDs, evidence and risk states. & Pass only if claims are supported, citations valid, risk constraints obeyed, and refusal boundaries respected; request repair for citation, grounding, wording, or boundary defects; refuse unrepairable defects. & \texttt{pass}, \texttt{repair\_required}, or \texttt{refuse\_required}; \texttt{refuse\_required} maps to \texttt{refused}. & \texttt{audit\_result} \\
\bottomrule
\end{tabular}
\end{table}

\subsection{Evidence-Constrained Generator}

If the Pre-Gate Request Filter allows the request and the Evidence Gate and Risk Gate do not require immediate refusal, the Evidence-Constrained Generator drafts a response under explicit evidence and risk constraints. The Generator may use only allowed evidence IDs and must avoid unsupported claims, invented pages, DOI-like references, or source attributions. Its output is a draft state, not a final answer.

The Generator receives \(q\), the evidence bundle, the Evidence Gate state, and the Risk Gate state. It outputs \texttt{generated\_response}, including cited evidence IDs and any safety de-escalation text required by the Risk Gate. The draft is always sent to the Independent Grounding Auditor.

\subsection{Bounded Citation/Grounding Repair}

Bounded Citation/Grounding Repair is deliberately narrow. It may repair only four defect classes: missing or invalid citations, grounding mismatches, unsupported wording, and missing refusal-boundary language. It may not repair an unsafe request into a safe one, add new medical or coaching recommendations, broaden the answer beyond the allowed evidence, or fill evidence gaps with model knowledge.

Repair receives the audit finding, the draft answer, and the allowed evidence bundle. It outputs \texttt{not\_needed}, \texttt{repaired}, or \texttt{repair\_failed}. A repaired draft is sent back to the Independent Grounding Auditor for re-check. If repair fails, the refusal decision is attributed to the Independent Grounding Auditor rather than to the Repair module itself.

\subsection{Trace Schema and Ablation Variants}

Each run stores a structured trace for every question. Table~\ref{tab:trace-schema} lists the minimal fields used to join benchmark metadata, evidence, outputs, and review notes; the same table also defines the ablation variants, each of which removes exactly one workflow function while preserving the remaining inputs, evidence modes, benchmark items, and final-status labels.

\begin{table}[t]
\centering
\scriptsize
\setlength{\tabcolsep}{2pt}
\caption{Compact trace schema and ablation variant definitions.}
\label{tab:trace-schema}
\label{tab:ablation-variants}
\begin{tabular}{@{}L{0.24\linewidth}L{0.66\linewidth}@{}}
\toprule
Item & Diagnostic definition \\
\midrule
\multicolumn{2}{@{}l}{\textit{Trace fields}} \\
\texttt{qid} & Stable benchmark item identifier. \\
\texttt{retrieved\_contexts} & Retrieval availability and noise, recorded as context counts and chunk IDs. \\
\texttt{evidence\_gate} / \texttt{risk\_gate} & Evidence sufficiency and safety-boundary constraints before generation. \\
\texttt{audit} / \texttt{repair} & Independent grounding/risk audit and bounded citation/grounding repair result. \\
\texttt{final\_status} & Released state: \texttt{answered}, \texttt{partial\_answer}, or \texttt{refused}. \\
\midrule
\multicolumn{2}{@{}l}{\textit{Ablation variants}} \\
\texttt{full} & Preserves all gates, generator, auditor, repair, and final-status labels. \\
\texttt{no\_gate} & Removes Evidence Gate refusal and partial/answerable gating; preserves retrieval, generation, audit, repair labels, and benchmark inputs. \\
\texttt{no\_audit} & Removes independent post-generation grounding, citation, and risk audit; preserves evidence/risk gates, generation, and final-status labels. \\
\texttt{no\_repair} & Removes Bounded Citation/Grounding Repair after audit; preserves gates, generator, audit findings, and refusal labels. \\
\bottomrule
\end{tabular}
\end{table}
